%! Rational Functions and Holes — Unit 1, Lesson 1.10 — Student Packet Cover Sheet
\documentclass[10pt]{article}
\usepackage{precalculus-article}
\usepackage{precalculus-boxes}
\usepackage{ltablex}
\keepXColumns

\begin{document}

% Full-bleed navy banner
\begin{tikzpicture}[remember picture, overlay]
  \fill[navy] (current page.north west)
    rectangle ([yshift=-0.9in]current page.north east);
\end{tikzpicture}
\vspace{-0.8in}

\begin{center}
  {\color{white}\textbf{\LARGE Precalculus}}\\[4pt]
  {\color{skymid}\large\textbf{Unit 1: Polynomial and Rational Functions}}\\[6pt]
  {\color{white}\textbf{\large Lesson 1.10 \quad Rational Functions and Holes}}
\end{center}

\vspace{0.22in}
\namedateperiod
\vspace{0.18in}

\begin{learningtargetbox}
\small
\begin{enumerate}[leftmargin=1.5em, itemsep=5pt]
  \item I can determine whether a rational function has a hole or a vertical asymptote at a
        shared zero of the numerator and denominator by comparing multiplicities.
  \item I can find the coordinates of a hole by canceling common factors and evaluating the
        simplified function at the shared zero, and express the result using limit notation.
\end{enumerate}
\end{learningtargetbox}

\vspace{0.14in}

\begin{tocbox}
\small
\renewcommand{\arraystretch}{1.5}
\begin{tabularx}{\linewidth}{@{} c l X r @{}}
\rowcolor{sky}
\textbf{\#} & \textbf{Component} & \textbf{Description} & \textbf{Score} \\
\midrule
1 & Warm-Up        & Spiral review: vertical asymptotes and factoring             & \blank{1.2cm} \\
2 & Group Activity & Missing Data Points --- classifying holes and VAs            & \blank{1.2cm} \\
3 & Exit Ticket    & Identify holes and write limit notation independently         & \blank{1.2cm} \\
4 & Homework       & Practice with holes, VAs, and limit notation                 & \blank{1.2cm} \\
\midrule
  & \multicolumn{2}{r}{\textbf{Total}} & \blank{1.2cm} \\
\end{tabularx}
\end{tocbox}

\end{document}
